\SourceLeaf{293}%
tet, zureichender Schärfe, so ist zu einer directen strengen Auflösung folgende Umformung am vortheilhaftesten, wobei die beiden ersten Formeln dieselben bleiben wie in der ersten Methode.

Vierte Methode
\[
\begin{aligned}
\operatorname{tang}s&=\cos T\operatorname{tang}r\\
\operatorname{tang}\lambda&=\frac{\operatorname{tang}T\sin s}{\cos(S-s)}\\
\operatorname{tang}t&=\sin T\sin r\operatorname{tang}(S-s)\\
\sin\tau&=\sin T\operatorname{tang}\tfrac{1}{2}r\sin s\\
\sin\sigma&=\operatorname{tang}t\operatorname{tang}\tfrac{1}{2}\lambda\cos(S-s)\\
S'&=S-s-\sigma\\
T'&=T-t-\tau
\end{aligned}
\]
\indent Diese vierte Methode lässt für die Schärfe nichts zu wünschen übrig; aber die unmittelbar in dieser Form geführte Rechnung erfordert ein etwas beschwerliches Interpoliren bei Bestimmung der kleinen Bögen durch die Logarithmen der Tangenten oder Sinus; man kann jedoch diesem Übelstande leicht ausweichen, indem man die trigonometrischen Functionen in Reihen entwickelt, wodurch man in den Stand gesetzt wird, ohne Nachtheil für die Schärfe, die Rechnungen vermittelst der Logarithmen der Zahlen zu führen. Es wird zureichend sein, von dieser Verwandlung nur die Hauptmomente hieher zu setzen.

Es sei
\[
\begin{aligned}
r\cos T&=s^0\\
r\sin T&=v
\end{aligned}
\]
Es wird dann, wenn zur Abkürzung die Grösse des Bogens von einer Secunde in Theilen des Halbmessers oder der Bruch \(\tfrac{\pi}{648000}\) durch \(\rho\) bezeichnet und \(r\) wie eine Grösse erster Ordnung betrachtet wird, bis auf Grössen fünfter Ordnung (ausschliesslich) genau
\[
s=s^0(1+\tfrac{1}{3}\rho\rho rr-\tfrac{1}{3}\rho\rho s^0s^0)=s^0(1+\tfrac{1}{3}\rho\rho vv)
\]
Setzt man dann ferner
\[
\begin{aligned}
v\operatorname{tang}(S-s)&=t^0\\
\frac{v}{\cos(S-s)}&=\lambda^0
\end{aligned}
\]
so wird
